\clearpage
\section{Introduction and Problem Statement}

\subsection{The Challenge: Stability and Functionality}

Materials discovery is a high-stakes filtration process: hundreds of thousands of candidate
compositions are screened to identify the few that can actually be synthesized and that have
the property an application needs. Early success in this pipeline hinges on overcoming two
critical challenges:

\begin{itemize}
  \item \textbf{Predicting Stability:} a material must be thermodynamically stable (on or near
  the convex hull of formation energies) to be made in the laboratory. Compositions resembling
  known stable compounds are statistically more likely to succeed, offering an early signal of
  viability.
  \item \textbf{Predicting Functionality:} electronic and magnetic behaviour (metal or
  insulator, magnetic or not, band gap inside the photovoltaic window) decides whether a
  material is useful for batteries, magnets, electronics or solar cells. Detecting these
  properties early avoids costly quantum-mechanical simulations and failed syntheses, saving
  both time and resources.
\end{itemize}

The essential question is clear: \emph{``Is this composition stable enough to be made, and does
it have the property we are looking for?''}

\subsection{The Data Dilemma}

Applying AI to this challenge is difficult due to an imbalance in available data:

\begin{enumerate}
  \item \textbf{Labeled Data is Scarce:} a reliable property label requires a density
  functional theory (DFT) calculation, costing hours to days of computing per material, or an
  even more expensive laboratory measurement. Consequently, high-quality labeled datasets
  remain limited compared with the space of possible compositions.
  \item \textbf{Unlabeled Data is Abundant:} high-throughput and generative searches such as
  Google DeepMind's GNoME propose hundreds of thousands of new compositions, most of which lack
  the property we want to predict.
\end{enumerate}

This raises a crucial question:

\begin{center}
\textbf{How can we harness the vast sea of unlabeled compositions to improve property
predictions, even when high-quality labeled data is scarce?}
\end{center}

Our solution is \textbf{Semi-Supervised Learning (SSL)}, which bridges the gap between scarce
labels and abundant unlabeled data. By leveraging unlabeled compositions intelligently, SSL aims
to enhance model performance and deliver more accurate and reliable property predictions.

\subsection{Complementary Approaches: SSL and Traditional Materials Informatics}

In addition to Semi-Supervised Learning, we implemented a traditional materials-informatics
pipeline for \textbf{photovoltaic band-gap classification}. This classical approach emphasizes
interpretability, descriptor optimization via genetic algorithms, and an OECD-style
applicability-domain assessment. We complete the study with \textbf{unsupervised learning}
(association-rule mining and clustering) and \textbf{self-supervised pre-training} on the
unlabeled pool. Together, these approaches provide complementary strategies: SSL leverages
unlabeled data for general property endpoints, unsupervised learning reveals the hidden
chemical structure of the data, and the traditional pipeline offers explainable, optimized
models for a well-characterized endpoint.
